\documentclass[12pt]{article}
\textwidth 7.0in
\oddsidemargin -0.4in
\evensidemargin -0.4in
\textheight 9.0in 
\pagestyle{empty}
\usepackage{amsmath, amssymb, amscd, amsfonts}
\usepackage{amsthm}
\usepackage{latexsym, epsfig, color, verbatim, xcolor, color}
\usepackage{url}

\usepackage{enumerate}
\begin{document}
% 8.5in paper width -2x1in margin = 6.5in text width
\setlength{\topmargin}{-2mm}



\newcommand{\sump}{\sideset{}{'}\sum}
\newcommand{\Aa}{{\alpha}}
\newcommand{\cD}{\mathcal{D}}
\newcommand{\C}{\mathbb{C}}
\newcommand{\N}{{\rm{N}}}
\newcommand{\re}{{\rm{max\,red}}}
\newcommand{\nonmax}{{\rm{nonmax}}}
\newcommand{\ma}{{\rm{max}}}
\newcommand{\PGL}{{\rm{PGL}}}
\newcommand{\diag}{{\rm{diag}}}
\newcommand{\sqf}{{\rm{sqf}}}
\newcommand{\Z}{\mathbb{Z}}
\renewcommand{\P}{\mathbb{P}}
\newcommand{\Q}{\mathbb{Q}}
\newcommand{\R}{\mathbb{R}}
\newcommand{\F}{\mathbb{F}}
\newcommand{\FF}{\mathcal{F}}
\newcommand{\OO}{\mathcal{O}}
\newcommand{\A}{\mathbb{A}}
\newcommand{\QQ}{\mathcal{Q}}
\newcommand{\calA}{\mathcal{A}}
\newcommand{\calS}{\mathcal{S}}
\newcommand{\calB}{\mathcal{B}}
\newcommand{\calP}{\mathcal{P}}
\newcommand{\calF}{\mathcal{F}}
\newcommand{\calZ}{\mathcal{Z}}
\newcommand{\calT}{\mathcal{T}}
\newcommand{\calO}{\mathcal{O}}
\newcommand{\Y}{\mathcal{Y}}
\newcommand{\GQ}{G_{\Q}}
\newcommand{\Qbar}{\bar{\Q}}
\newcommand{\ra}{\rightarrow}
\newcommand\Hom{\operatorname{Hom}}
\newcommand\Tr{\operatorname{Tr}}
\newcommand\Aut{\operatorname{Aut}}
\newcommand\Sym{\operatorname{Sym}}
\newcommand\Disc{\operatorname{Disc}}
\newcommand\Vol{\operatorname{Vol}}
\newcommand\Out{\operatorname{Out}}
\newcommand\im{\operatorname{im}}
\renewcommand\div{\operatorname{div}}
\newcommand\irr{\textnormal{irr}}
\newcommand\Stab{\operatorname{Stab}}
\newcommand\Gal{\operatorname{Gal}}
\newcommand\Res{\operatorname{Res}}
\newcommand\SL{\operatorname{SL}}
\newcommand\GL{\operatorname{GL}}
\newcommand\SO{\operatorname{SO}}
\newcommand\ur{\operatorname{ur}}
\newcommand\sr{\operatorname{sr}}
\newcommand\htt{\operatorname{ht}}
\newcommand{\cO}{\mathcal{O}}
\newcommand{\co}{\mathcal{O}}
\newcommand\codim{\operatorname{codim}}


\newcommand{\odr}{\co_{D1^2}}		% Ramify
\newcommand{\ods}{\co_{D11}}		% Split
\newcommand{\odi}{\co_{D2}}		% Inert
\newcommand{\odg}{\co_{D\rm{ns}}}	% Generic
% Common component
\newcommand{\ocs}{\co_{C\rm{s}}}	% Singular
\newcommand{\ocg}{\co_{C\rm{ns}}}	% Generic
% essentially in Two variables
\newcommand{\ots}{\co_{B11}}		% Split
\newcommand{\oti}{\co_{B2}}		% Inert

% 11in paper height -2x1in margin = 9in text height


\begin{center}{\bf Homework 1, Math 701 -- Frank Thorne (thorne@math.sc.edu)}
\end{center}
\begin{center}
{\bf Instructions:} You are welcome and encouraged to collaborate, but please write up your own solutions. \\

Due Wednesday, September 9.
\end{center}

\begin{enumerate}
%
%\item
%Let $V$ be a vector space, possibly infinite dimensional.
%\begin{enumerate}
%\item
%Prove that the map
%\[
%v \rightarrow \{ f \rightarrow f(v) \}
%\]
%defines an injective map from $V$ to the double dual $V^{**} = (V^*)^*$.
%\item
%Prove that this map is not surjective if $V$ is infinite dimensional.
%\end{enumerate}
%
%\item
%\begin{enumerate}
%\item
%Exhibit a $2 \times 2$ matrix with real entries which does not have any eigenvalues over $\mathbb{R}$. 
%\item 
%Characterize, as explicitly as you can, the set of all such matrices,
%\end{enumerate}
%
%\item 
%Suppose that $\phi$ is a linear transformation whose Jordan canonical form $M$ consists of blocks of size $k_i$ and with diagonal entries $\lambda_i$, for $i = 1, ..., r$.
%
%\begin{itemize}
%\item Write down explicitly what that means. (Here is probably where you want to read a more detailed exposition of Jordan form in your favorite book.)
%\item Prove that each $\lambda_i$ is an eigenvalue of $\phi$, and describe all of the associated eigenvectors in terms of $M$. 
%\item Prove that no scalar other than the $\lambda_i$ is an eigenvalue of $\phi$.
%\item Let $f(x) := (x- \lambda_i)^{k_i}$. Explain why this is the characteristic polynomial of $\phi$. 
%\item The {\bf Cayley-Hamilton theorem} asserts that $f(\phi) = 0$. Prove it.
%\end{itemize}
%
%\item 
%Prove that the Jordan canonical form is not uniquely determined, by showing that the matrices
%\[
%\begin{pmatrix} 2 & 1 & 0 \\ 0 & 2 & 0 \\ 0 & 0 & 3 \end{pmatrix}, \ \ \ \ 
%\begin{pmatrix} 3 & 0 & 0 \\ 0 & 2 & 0 \\ 0 & 1 & 2 \end{pmatrix}
%\]
%are conjugate. 
%
%\item
%Suppose that $\phi$ is a linear transformation with characteristic polynomial $f(x) = (x - 2)^2 (x - 3)^2$.
%
%Determine how many possible Jordan forms there are for $\phi$, and write them all down up to conjugacy. 
%(In other words only write down a single Jordan form, when there are multiple choices and they are conjugate to each other.)
%Only write down a single Jordan form for 
%
%
%
%Better still:
%can you characterize all such matrices?)
%
%
%Give a proof that every basis of $\mathbb{R}^2$ has exactly two vectors. {\bf Your proof should be as minimal as possible} -- don't use {\bf any} lemmas
%or propositions described in class unless you prove them.

\item 
The {\itshape Fibonacci numbers} are defined by $F_1 = 1$, $F_2 = 1$, and $F_{n + 1} = F_n + F_{n - 1}$ for $n \geq 2$.
Note that they satisfy
\[
M
\begin{bmatrix}
F_n \\ F_{n - 1}  
\end{bmatrix}
=
\begin{bmatrix}
F_{n + 1} \\ F_{n}  
\end{bmatrix}
\]
where $M := \begin{bmatrix}
1 & 1 \\ 1 & 0 
\end{bmatrix}$, for each $n \geq 2$.
\begin{enumerate}[(a)]
\item Compute the characteristic polynomial, eigenvalues, and a basis of eigenvectors for $M$. 
\item Compute a diagonal matrix $D$ and an invertible matrix $P$ for which $D = P^{-1} M P$. Interpreting $P$ as a change 
of basis matrix, describe the linear transformations and bases represented by $D$, $P$, and $M$.
\item Sketch a proof that the set of functions $f : \Z^+ \rightarrow \R$ satisfying $f(n + 1) = f(n) + f(n - 1)$ for $n \geq 2$
forms a two-dimensional real vector space $V$. 
\item Using your work above, find a basis for $V$ consisting of functions satisfying $f(n) = f(1) \lambda^{n - 1}$ for some
$\lambda \in \R$. 
\item Write $F$ in terms of this basis, thereby computing an explicit closed formula for the Fibonacci numbers.
\end{enumerate}

\item 
Let $V$ be the vector space of {\itshape affine plane conics}
\[
ax^2 + bxy + cy^2 + dx + ey + f = 0
\]
with coefficients in $\mathbb{C}$. Then $V$ is a vector space of dimension $6$. 

\begin{enumerate}[(a)]
\item Fix $(x_0, y_0)$, and let $W \subseteq V$ be the subspace of conics $g$ for which $g(x_0, y_0) = 0$. Explain briefly why $W$
is a subspace, and why it has dimension $5$ for any choice of $(x_0, y_0)$. 
\item Now, for $k \geq 1$, consider the following statement:
{\itshape For any $k$ distinct points $(x_1, y_1), \dots, (x_k, y_k)$, let $W_k \subseteq V$ be the subspace of conics $g$ for which $g(x_i, y_i) = 0$
for each $i$. Then $W_k$ has dimension $6 - i$.}

For which $k \geq 1$ is the above statement true or false? (Prove your assertion.)

\item A conic $g(x, y)$ is {\itshape singular} if it contains some point $(x_0, y_0)$ at which both partial derivatives vanish. Find (with proof) an example
of a nonsingular conic, and a nonzero example of a singular conic.

\item Find a basis for $V$ consisting of nonsingular conics. 
\end{enumerate}

\item
Let $V$ be the real vector space consisting of all continuous functions from $\mathbb{R}$ to itself. Prove that $V$ is not finite dimensional.

\item (Lagrange Interpolation, following Friedberg, Insel, and Spence)
Let $V$ be the $n + 1$-dimensional vector space of polynomials of degree $\leq n$, over a field $F$, and let $c_0, \dots, c_n$ be
distinct scalars in $F$.
\begin{enumerate}[(a)]

\item 
For $0 \leq i \leq n$, define $f_i \in V^*$ by $f_i(p) = p(c_i)$. Prove that the $f_i$ form a basis for $V^*$.
\item 
Prove that there exist unique polynomials $p_0, \dots, p_n$ with $p_i(c_j) = \delta_{ij}$.
\item 
For any scalars $a_0, \dots, a_n$ (not necessarily distinct), prove that there exists a unique polynomial $q$ of degree $\leq n$,
such that $q(c_i) = a_i$ for each $i$. Determine a formula for $q$ in terms of the above.

\end{enumerate}

\item (Suggested by ChatGPT)
For each $n$, let $P_n$ be the vector space of polynomials over $\R$ of degree at most $n$.
\begin{enumerate}[(a)]
\item 
Let $D : P_3 \mapsto P_2$ be the differentiation map. Briefly explain why $D$ is a homomorphism. Compute a matrix representing $D$
(with respect to bases you should choose and specify).
\item 
Similarly to the previous question, for $a \in \R$ define $\textnormal{ev}_a$ to be the `evaluation at $a$ map', that is 
$\textnormal{ev}_a(p) = p(a)$. Explain why this is an element of $P_n^*$, and describe $D^*(\textnormal{ev}_a)$.
\item 
Define $I : P_n \mapsto \R$ by
\[
I(p) = \int_0^1 p(x) dx.
\]
Explain why $I \in P_n^*$. 
\item
Compute $D^*(I)$.
\end{enumerate}

\end{enumerate}

\end{document}


\item
({\bf Row Rank Equals Column Rank})
The {\bf row rank} and {\bf column rank} of an $n \times m$ 
matrix $A$ are the dimensions of the subspaces of $F^n$
and $F^m$, respectively, spanned by its rows and columns respectively.

Prove\footnote{Credit: I learned of this proof from Marc van Leeuwen} that the row
rank and column rank are equal as follows:
\begin{itemize}
\item
Show that the column rank is the smallest number $r$ for which the homomorphism
$\phi$ corresponding to $A$ can be written as a composition of two linear maps
$F^m \rightarrow F^r \rightarrow F^n$.
\item
Conclude that $r$ is the smallest number for which one can write $A = BC$,
with $B$ an $n \times r$ and $C$ an $r \times m$ matrix respectively.
\item
Taking transposes, conclude the result.
\end{itemize}
\item
Let $V$ be a vector space, possibly infinite dimensional.
\begin{enumerate}
\item
Prove that the map
\[
v \rightarrow \{ f \rightarrow f(v) \}
\]
defines an injective map from $V$ to the double dual $V^{**} = (V^*)^*$.
\item
Prove that this map is not surjective if $V$ is infinite dimensional.
\end{enumerate}

\item
Let $V$ be a finite dimensional vector space, and let $A$ and $B$ be the matrices of a linear transformation $\phi$
with respect to two different bases.

Prove that $B = M A M^{-1}$ for some matrix $M$, and describe $M$ explicitly.

\item
Exhibit a $2 \times 2$ matrix with real entries which does not have any eigenvalues over $\mathbb{R}$. (Better still:
can you characterize all such matrices?)

\end{enumerate}

\end{document}
